\chapter{Haar Measure, Duality and the \texorpdfstring{$p$}{p}-adic Numbers}

\begin{center}\itshape
Lebesgue measure is the measure that does not care where you are. Every reasonable group has such a measure, and its characters
form a second group that mirrors the first. On the real line, the mirror is the line itself. On the $p$-adic numbers, the mirror is the field itself.
\end{center}

In 1924 Stefan Banach and Alfred Tarski cut a solid ball into finitely many pieces and reassembled them, by rotations and translations, into
two balls of the same size as the first. Nothing is wrong with the logic: the pieces are so wild that they have no volume, and the
axiom of choice builds them. But a question remained, which Volume~7 (Chapter~7) left open. In the plane, Banach had shown in 1923
that nothing of the kind can happen: there is a notion of area, finitely additive and invariant under all rigid motions, defined for
\emph{every} bounded set. Why can one cut and reassemble balls in space but not discs in the plane? John von Neumann answered in 1929, and his
answer is not about space but about the \emph{group} of motions: in dimension three it contains a free group, which can be
duplicated by its own structure, while in dimension two it is built from commutative groups, which admit invariant averages. He
called such groups \emph{messbar}, measurable; Mahlon Day later renamed them \emph{amenable}, with a pun on ``mean''.

Invariant averages are the theme of this chapter. On a compact group they are integrals against a probability measure that every
translation preserves, the \emph{Haar measure}, found by Alfr\'ed Haar in 1933 for a large class of groups and by Andr\'e Weil in general. Haar measure is what
makes Fourier analysis possible on a group: Chapter~18 used Lebesgue measure on $\R^d$, Chapter~20 counting measure on finite groups, and
the same theory runs on the circle, on the integers, and on the $p$-adic numbers. The field $\Q_p$ is constructed in Volume~13. This chapter gives that field a Haar measure and proves that it is self-dual: the new page of Pontryagin's dictionary. The ring that holds $\R$ and every $\Q_p$ at once, the adèles, is assembled in Volume~13, and the Poisson formula on it is the functional equation of $\zeta$. What that argument uses from here is the measure and the duality.

\section{Topological groups and Haar measure}

A \emph{topological group} is a group with a topology for which $(x,y)\mapsto xy$ and $x\mapsto x^{-1}$ are continuous. To keep the measure theory
of Volume~7 at hand, we assume throughout that our groups are \emph{metrizable}, \emph{locally compact} (every point has a compact neighbourhood) and
\emph{$\sigma$-compact} (a countable union of compact sets). Examples: $\R^d$; the circle $\mathbb T=\{z\in\C:|z|=1\}$; every countable group with the discrete
topology; $GL_n(\R)$, an open subset of $\R^{n^2}$; and, below, $\Q_p$ and $\Z_p$.

\begin{definition}
A \emph{left Haar measure} on $G$ is a Borel measure $\mu$ that is finite on compact sets, positive on non-empty open sets, and \emph{left invariant}:
$\mu(gE)=\mu(E)$ for all $g\in G$ and Borel $E$. Right Haar measures are defined with $Eg$. On an abelian group, written additively, we simply
say \emph{Haar measure}.
\end{definition}

Lebesgue measure on $\R^d$, arc length on $\mathbb T$ and counting measure on a discrete group are Haar measures. Products of Haar measures are
Haar measures on products of groups. Here is a non-commutative example.

\begin{proposition}\label{prop:glnhaar}
The measure on $GL_n(\R)$ defined by
\[\mu(E)=\int_E|\det X|^{-n}\dd X,\]
where $\dd X$ is Lebesgue measure on $\R^{n^2}$, is a left and right Haar measure.
\end{proposition}
\begin{proof}
For $A\in GL_n(\R)$, the map $L_A:X\mapsto AX$ is linear on $\R^{n^2}$; it acts on each of the $n$ columns of $X$ by $A$, so $\det L_A=(\det A)^n$. By the change of
variables formula (Volume~7, Theorem~9.8),
\[\mu(AE)=\int_{AE}|\det X|^{-n}\dd X=\int_E|\det AY|^{-n}|\det A|^n\dd Y=\mu(E).\]
Right multiplication acts on the rows and is handled in the same way. The density is continuous and positive on the open set $GL_n(\R)$, so $\mu$ is
finite on compact sets and positive on open sets.
\end{proof}

Left and right Haar measures need not agree. On the group of affine maps $x\mapsto ax+b$ ($a>0$), the left Haar measure is $\frac{\dd a\dd b}{a^2}$ and the
right one is $\frac{\dd a\dd b}a$ (Exercise~\ref{ex:axb}). In general, $E\mapsto\mu(Eg)$ is again a left Haar measure, so by uniqueness it is $\Delta(g)\mu$ for a
number $\Delta(g)>0$, the \emph{modular function}; groups with $\Delta=1$ are called \emph{unimodular}. Abelian, compact and discrete groups are unimodular.

\begin{theorem}[uniqueness]\label{thm:haaruniq}
Two Haar measures on an abelian group are proportional.
\end{theorem}
\begin{proof}
Let $\mu$ and $\nu$ be Haar measures on $G$, and let $f,g$ be continuous with compact support; write $\check g(x)=g(-x)$. All the functions below are continuous with compact
support in $G\times G$, and $\mu$, $\nu$ are $\sigma$-finite, so Fubini's theorem (Volume~7, Theorem~9.6) applies. Using the invariance of $\mu$ ($x\to x+y$),
Fubini, the invariance of $\nu$ ($y\to y-x$), Fubini again, and the invariance of $\mu$ ($x\to x+y$),
\[\begin{aligned}\int f\dd\mu\int g\dd\nu&=\int\!\!\int f(x+y)g(y)\dd\mu(x)\dd\nu(y)\\&=\int\!\!\int f(x+y)g(y)\dd\nu(y)\dd\mu(x)\\&=\int\!\!\int f(y)g(y-x)\dd\nu(y)\dd\mu(x)\\&=\int f(y)\int g(y-x)\dd\mu(x)\dd\nu(y)=\int f\dd\nu\int\check g\dd\mu.\end{aligned}\]
Fix $g\ge0$, not identically $0$; then $\int g\dd\nu>0$, since $\nu$ charges the open set $\{g>0\}$. So $\int f\dd\mu=c\int f\dd\nu$ for every $f$, with
$c=\int\check g\dd\mu/\int g\dd\nu$. It remains to see that a Borel measure finite on compact sets is determined by these integrals. If $K$ is compact and $V$ a
compact neighbourhood of $K$, the functions $f_n(x)=\max\bigl(0,1-n\,d(x,K)\bigr)$ are continuous, supported in $V$ for $n$ large, and decrease to $\mathbf 1_K$; by
dominated convergence $\mu(K)=c\nu(K)$. Compact sets form a $\pi$-system; they generate the Borel sets, since every closed set $F$ is the countable union
of the compact sets $F\cap K_n$, where $G=\bigcup K_n$ with $K_n$ compact and increasing. Theorem~6.5 of Volume~7 concludes.
\end{proof}

The same argument, with more care about left and right, proves uniqueness for every locally compact group. Existence is harder.

\begin{remark}
\emph{Haar's theorem} (Haar 1933 for compact metrizable groups and separable locally compact groups; Weil 1940 in general) says that every locally compact group
has a left Haar measure. We shall not need it in full generality: the groups of this chapter either have explicit Haar measures, or are compact and
abelian, and the next section proves existence for those.
\end{remark}

\section{Invariant means and amenability}

Let $G$ be a group with the discrete topology, and $\ell^\infty(G)$ the bounded real functions on $G$. A \emph{mean} is a linear form $m$ on $\ell^\infty(G)$ that is
positive ($m(f)\ge0$ when $f\ge0$) with $m(\mathbf 1)=1$; it is \emph{invariant} if $m(\lambda_gf)=m(f)$ for all $g$, where $\lambda_gf(x)=f(g^{-1}x)$. The group is
\emph{amenable} if it has an invariant mean. A Banach limit (Exercise~\ref{ex:banachlimit}) is an invariant mean on the semigroup $\N$; on a finite group, the average
is one. Invariant means on infinite groups cannot be written down: they come from a fixed point theorem.

\begin{theorem}[Markov and Kakutani, 1936--38]\label{thm:markovkakutani}
Let $E$ be a Banach space and $K\subset E^*$ non-empty, convex and weak* compact. Let $\mathcal T$ be a family of maps $K\to K$ that are affine (they preserve
convex combinations), weak* continuous, and commute with each other. Then the maps of $\mathcal T$ have a common fixed point.
\end{theorem}
\begin{proof}
For $T\in\mathcal T$ and $n\ge1$ let $A^T_n=\frac1n(I+T+\cdots+T^{n-1})$, an affine continuous map $K\to K$ ($K$ is convex). These maps commute with each other: since the $T$
are affine, $A^S_nA^T_m$ and $A^T_mA^S_n$ both send $x$ to $\frac1{nm}\sum_{i<n,j<m}S^iT^jx$, and $S^iT^j=T^jS^i$. The sets $A^T_n(K)$ are weak* compact. Any finite
family of them has a common point: if $B$ is the composition of $A^{T_1}_{n_1},\dots,A^{T_k}_{n_k}$, then for each $i$, $B=A^{T_i}_{n_i}\circ C_i$ with $C_i(K)\subset K$, so
$B(K)\subset A^{T_i}_{n_i}(K)$. By compactness, there is $x_0$ in all the sets $A^T_n(K)$.

Fix $T$ and $n$, and write $x_0=A^T_ny$ with $y\in K$. Since $T$ is affine, $Tx_0-x_0=\frac1n(T^ny-y)$. For $e\in E$, the weak* continuous function $z\mapsto\langle z,e\rangle$ is
bounded on the compact set $K$, by some $M_e$, so $|\langle Tx_0-x_0,e\rangle|\le\frac{2M_e}n$. Letting $n\to\infty$, $\langle Tx_0,e\rangle=\langle x_0,e\rangle$ for every $e$: $Tx_0=x_0$.
\end{proof}

\begin{corollary}\label{cor:abelianamenable}
\begin{enumerate}
\item[(a)] Every abelian group is amenable.
\item[(b)] Every compact metrizable abelian group has a Haar probability measure, and it is unique.
\end{enumerate}
\end{corollary}
\begin{proof}
(a) Let $K$ be the set of means, inside the dual of the Banach space $\ell^\infty(G)$. A mean has norm $1$ (for $|f|\le\|f\|_\infty\mathbf 1$, positivity gives
$|m(f)|\le\|f\|_\infty$), and the conditions defining a mean are weak* closed, so $K$ is weak* compact by Banach--Alaoglu (Theorem~\ref{thm:alaoglu}). It is convex and
non-empty (evaluation at a point is a mean). The maps $T_gm=m\circ\lambda_g$ are affine, weak* continuous and send means to means; they commute because $G$ is
abelian. A common fixed point is an invariant mean.

(b) By the Riesz--Markov theorem (Theorem~\ref{thm:rieszmarkov}), the probability measures on $G$ form a weak* compact convex subset $K$ of $C(G)^*$: the positive
forms with $\varphi(\mathbf 1)=1$. Translations $T_g\mu=\mu(\cdot-g)$, that is $\int f\dd(T_g\mu)=\int f(x+g)\dd\mu(x)$, are affine, weak* continuous and commute. A fixed point is a
translation invariant probability $\mu$. If $\mu(U)=0$ for a non-empty open $U$, finitely many translates of $U$ cover the compact group, and $\mu(G)=0$, absurd.
So $\mu$ is a Haar measure, and Theorem~\ref{thm:haaruniq} gives uniqueness among probabilities.
\end{proof}

\begin{proposition}\label{prop:amenableext}
Let $N$ be a normal subgroup of $G$. If $N$ and $G/N$ are amenable, so is $G$. In particular, finite extensions of abelian groups, and extensions of abelian
groups by abelian groups, are amenable.
\end{proposition}
\begin{proof}
Let $m_N$ and $m_Q$ be invariant means on $N$ and on $Q=G/N$. For $f\in\ell^\infty(G)$ put $F(gN)=m_N\bigl(n\mapsto f(gn)\bigr)$. This only depends on the coset: replacing $g$ by
$gn_0$ replaces the function $n\mapsto f(gn)$ by its translate $n\mapsto f(gn_0n)$, which has the same mean by invariance (with $\lambda_{n_0^{-1}}$). Set $m(f)=m_Q(F)$. It is a
mean, and for $h\in G$, $\lambda_hf$ gives $F'(gN)=m_N(n\mapsto f(h^{-1}gn))=F(h^{-1}gN)$, that is $F'=\lambda_{hN}F$, so $m(\lambda_hf)=m(f)$. Finite groups are amenable
(average), and the last sentence follows with Corollary~\ref{cor:abelianamenable}.
\end{proof}

The group $\mathrm{Isom}(\R^2)$ of rigid motions of the plane is an extension of the group $O(2)$ of rotations and reflections, itself a
two-element extension of the abelian $SO(2)$, by the abelian group of translations $\R^2$. As a discrete group it is therefore amenable, and so is $\mathrm{Isom}(\R)$.

\begin{theorem}[Banach, 1923]\label{thm:banachmeasure}
For $d=1,2$ there is a function $\mu$ on all bounded subsets of $\R^d$, finitely additive, invariant under rigid motions, and equal to Lebesgue measure on
Lebesgue measurable sets.
\end{theorem}
\begin{proof}
Let $V$ be the space of bounded real functions with bounded support on $\R^d$, and $p(f)$ the \emph{upper integral} of $f\in V$: the infimum of $\int h\dd\lambda$ over the
measurable $h\in V$ with $h\ge f$; $p$ is sublinear, and $p(f)=\int f$ when $f$ is measurable. By Hahn--Banach (Theorem~\ref{thm:hb}) there is a linear $L$ on $V$ with $L\le p$ and $L(f)=\int f$
for measurable $f$. It is positive: if $f\ge0$ then $L(-f)\le p(-f)\le0$. Moreover $p$ is invariant under rigid motions, as $\lambda$ is.

For a bounded $A$, the function $\varphi_A(g)=L(\mathbf 1_{g^{-1}A})$ on $G=\mathrm{Isom}(\R^d)$ satisfies $0\le\varphi_A\le p(\mathbf 1_{g^{-1}A})=p(\mathbf 1_A)$, so it is bounded. Let $m$ be an invariant
mean on $G$ and $\mu(A)=m(\varphi_A)$. Linearity of $L$ and $m$ gives finite additivity. For $h\in G$, $\varphi_{hA}(g)=\varphi_A(h^{-1}g)=\lambda_h\varphi_A(g)$, so $\mu(hA)=\mu(A)$. If
$A$ is measurable, $\varphi_A$ is the constant $\lambda(A)$, so $\mu(A)=\lambda(A)$.
\end{proof}

In particular, no disc can be cut into finitely many pieces that rigid motions reassemble into two discs: $\mu$ would give $\pi r^2=2\pi r^2$. In space, the
rotations alone defeat every such $\mu$, because they contain a free group. Write $F_2$ for the free group on two letters $a,b$: its elements are the
\emph{reduced words} in $a^{\pm1},b^{\pm1}$ (no letter next to its inverse), multiplied by concatenation and cancellation. Let $W(s)$ be the set of reduced words
beginning with the letter $s$.

\begin{proposition}\label{prop:f2paradox}
$F_2=\{e\}\sqcup W(a)\sqcup W(a^{-1})\sqcup W(b)\sqcup W(b^{-1})$, while $F_2=W(a)\sqcup aW(a^{-1})=W(b)\sqcup bW(b^{-1})$. Consequently $F_2$ is not amenable.
\end{proposition}
\begin{proof}
The first decomposition sorts words by their first letter. If $w$ does not begin with $a$, then $a^{-1}w$ is reduced and begins with $a^{-1}$, so $w\in aW(a^{-1})$;
and a word $aw'$ with $w'\in W(a^{-1})$ cancels to a word not beginning with $a$. This proves the second decomposition, and the third likewise. If $m$ were an
invariant mean, writing $m(E)=m(\mathbf 1_E)$, then $1=m(W(a))+m(aW(a^{-1}))=m(W(a))+m(W(a^{-1}))$ and likewise for $b$, so $m(F_2)\ge2$, absurd.
\end{proof}

\begin{figure}[ht]
\centering
\fig{ch08-freegroup}
\caption{The Cayley graph of the free group $F_2$ to depth $3$, coloured by first letter: $W(a)$ to the right, $W(a^{-1})$ to the left, $W(b)$ up, $W(b^{-1})$ down.
Multiplying $W(a^{-1})$ by $a$ moves the whole left branch to the centre and turns it into everything except the right branch: one quarter of the tree becomes
three quarters.}
\label{fig:freegroup}
\end{figure}

\begin{lemma}[Hausdorff, 1914]\label{lem:freerot}
Let $\phi$ and $\rho$ be the rotations of $\R^3$ by the angle $\arccos\frac13$ about the $z$-axis and about the $x$-axis. Then no non-empty reduced word in
$\phi^{\pm1},\rho^{\pm1}$ is the identity: $\phi$ and $\rho$ generate a free group.
\end{lemma}
\begin{proof}
In matrix form,
\[\phi^{\pm1}=\frac13\begin{pmatrix}1&\mp2\sqrt2&0\\\pm2\sqrt2&1&0\\0&0&3\end{pmatrix},\qquad\rho^{\pm1}=\frac13\begin{pmatrix}3&0&0\\0&1&\mp2\sqrt2\\0&\pm2\sqrt2&1\end{pmatrix}.\]
We show: if $w$ is a reduced word of length $k\ge1$ ending (on the right) with $\phi^{\pm1}$, then $w(1,0,0)=3^{-k}(a,b\sqrt2,c)$ with $a,b,c\in\Z$ and $3\nmid b$. In particular
$w(1,0,0)\ne(1,0,0)$. For $k=1$, $\phi^{\pm1}(1,0,0)=\frac13(1,\pm2\sqrt2,0)$. Otherwise $w=\sigma w'$ with $w'(1,0,0)=3^{-(k-1)}(a',b'\sqrt2,c')$, and multiplying out,
\[\phi^{\pm1}:(a,b,c)=(a'\mp4b',\ b'\pm2a',\ 3c'),\qquad\rho^{\pm1}:(a,b,c)=(3a',\ b'\mp2c',\ c'\pm4b').\]
So $a,b,c$ are integers, and we check $3\nmid b$ according to the first two letters $\sigma\tau$ of $w$, writing $w'=\tau w''$ with the same notation $(a'',b'',c'')$
(for $k=2$, $w''$ is empty and $(a'',b'',c'')=(1,0,0)$). If $\sigma=\phi^{\pm1}$ and $\tau=\rho^{\pm1}$, then $a'=3a''$ and $b=b'\pm6a''\equiv b'$. If $\sigma=\rho^{\pm1}$ and $\tau=\phi^{\pm1}$,
then $c'=3c''$ and $b\equiv b'$. If $\sigma=\tau=\phi^{\pm1}$ (the same sign, as $w$ is reduced), then $\pm2a'=\pm2a''-8b''=(b'-b'')-8b''$, so $b=2b'-9b''\equiv2b'$; the case
$\sigma=\tau=\rho^{\pm1}$ is identical. In all cases $3\nmid b$ by induction.

Finally, if a non-empty reduced word $w$ ending with $\rho^{\pm1}$ were the identity, so would be $\phi w\phi^{-1}$, which after the possible cancellation of a leading
$\phi^{-1}$ in $w$ is a non-empty reduced word ending with $\phi^{-1}$: impossible.
\end{proof}

\begin{theorem}[Hausdorff, 1914]\label{thm:hausdorff}
There is no finitely additive, rotation invariant function $\mu$ on all subsets of the sphere $S^2$ that agrees with surface area on Borel sets.
\end{theorem}
\begin{proof}
Let $F$ be the group generated by $\phi$ and $\rho$, which is countable and free. Each rotation $w\ne e$ in $F$ fixes exactly two points of $S^2$; let $D$ be the
countable set of all these points. If $wx=x$ then $(gwg^{-1})(gx)=gx$, so $D$ and $X=S^2\setminus D$ are $F$-invariant, and $F$ acts on $X$ without fixed points. By the axiom
of choice, let $M\subset X$ contain exactly one point of each orbit. Every $x\in X$ is then $gm$ for a unique $g\in F$ and $m\in M$ (if $gm=g'm$ then $g^{-1}g'$ fixes $m$, so
$g=g'$). Hence, by Proposition~\ref{prop:f2paradox} transported to $F$,
\[\begin{aligned}X&=M\sqcup W(\phi)M\sqcup W(\phi^{-1})M\sqcup W(\rho)M\sqcup W(\rho^{-1})M,\\X&=W(\phi)M\sqcup\phi W(\phi^{-1})M=W(\rho)M\sqcup\rho W(\rho^{-1})M.\end{aligned}\]
If $\mu$ existed, then $\mu(D)=0$ (it is contained in Borel sets of arbitrarily small area) and $\mu(X)=4\pi$. Rotation invariance gives
$\mu(W(\phi)M)+\mu(W(\phi^{-1})M)=4\pi$, and the same with $\rho$, so $\mu(X)\ge8\pi$: a contradiction.
\end{proof}

\begin{remark}
Two more pages turn this into the Banach--Tarski paradox: the countable set $D$ can be absorbed by a rotation, the sphere becomes a ball minus its centre by
taking radii, and the centre is absorbed by a rotation about an axis off the centre. See G.~Tomkowicz and S.~Wagon, \emph{The Banach--Tarski Paradox}, Chapters~1--3.
The dichotomy is now clear: $\mathrm{Isom}(\R^d)$ is amenable for $d\le2$ and contains $F_2$ for $d\ge3$. Tarski proved in 1938 that a group acting on a set admits an
invariant finitely additive probability on all subsets if and only if the set has no paradoxical decomposition.
\end{remark}

\section{Characters and duality}

Let $G$ be a locally compact abelian group, written additively. A \emph{character} is a continuous homomorphism $\chi:G\to\mathbb T$; the characters form the
\emph{dual group} $\widehat G$ under pointwise multiplication. For finite groups, these are the characters of Chapter~20. We write $e(t)=e^{2\pi it}$.

\begin{theorem}[some dual groups]\label{thm:duals}
\begin{enumerate}
\item[(a)] Every character of $\R^d$ is $x\mapsto e(\xi\cdot x)$ for a unique $\xi\in\R^d$: $\widehat{\R^d}\cong\R^d$.
\item[(b)] Every character of $\Z$ is $n\mapsto z^n$ for a unique $z\in\mathbb T$: $\widehat{\Z}\cong\mathbb T$.
\item[(c)] Every character of $\mathbb T$ is $z\mapsto z^k$ for a unique $k\in\Z$: $\widehat{\mathbb T}\cong\Z$.
\end{enumerate}
\end{theorem}
\begin{proof}
(a) By Proposition~\ref{prop:dualprod}(a) (whose proof works verbatim for continuous characters) it suffices to treat $d=1$. Let $\chi$ be a character of $\R$. As $\chi(0)=1$
and $\chi$ is continuous, there is $h>0$ with $c=\int_0^h\chi(t)\dd t\ne0$. Then $c\chi(x)=\int_0^h\chi(x+t)\dd t=\int_x^{x+h}\chi(t)\dd t$, which is $C^1$ in $x$, so $\chi$ is $C^1$ and
$c\chi'(x)=\chi(x+h)-\chi(x)=\chi(x)(\chi(h)-1)$. So $\chi'=\kappa\chi$ for a constant $\kappa$, and $\chi(x)=e^{\kappa x}$; $|\chi|=1$ forces $\kappa=2\pi i\xi$ with $\xi$ real.
Uniqueness: $e(\xi x)=e(\xi'x)$ for all $x$ forces $\xi=\xi'$ (differentiate at $0$).

(b) A homomorphism on $\Z$ is determined by $z=\chi(1)$, which can be any point of $\mathbb T$.

(c) If $\chi$ is a character of $\mathbb T$, then $t\mapsto\chi(e(t))$ is a character of $\R$, so it is $e(\xi t)$ by (a), and it is $1$ at $t=1$, so $\xi=k\in\Z$.
\end{proof}

So $\R$ is its own dual, while $\Z$ and $\mathbb T$ are each other's duals. In every case the dual of the dual is the group we started from, as for finite groups
(Theorem~\ref{thm:finiteduality}), and compact groups have discrete duals and vice versa.

\begin{remark}
The \emph{Pontryagin--van Kampen duality theorem} (Pontryagin 1934 for compact and discrete groups, van Kampen 1935 in general) says that for every locally compact
abelian group $G$, with the topology of uniform convergence on compact sets on $\widehat G$, the map $G\to\widehat{\widehat G}$ is an isomorphism of topological groups. Its proof
needs Fourier analysis on $G$ itself, with the Haar measure: one defines $\hat f(\chi)=\int_Gf\overline\chi\dd\mu$, proves a Plancherel theorem for a suitably normalised Haar measure on
$\widehat G$, and deduces that characters separate points. Fourier series ($G=\mathbb T$), the Fourier transform of Chapter~18 ($G=\R^d$), the discrete transform of Chapter~20, and the
$p$-adic transform below are all instances; see Folland, \emph{A Course in Abstract Harmonic Analysis}, Chapter~13. We verify the theorem directly on the groups we need.
\end{remark}

\begin{beauty}
The circle and the integers are dual, and the dictionary between them is the theory of Fourier series: a function on the circle is a sum of characters $z^k$
indexed by integers, and a sequence of integers is ``a function on the circle'' through its generating series. Smoothness on one side becomes decay on the other,
convolution becomes multiplication, Parseval holds. Pontryagin's theorem says that this dictionary exists for every locally compact abelian group, and the
$p$-adic numbers below give a new, totally disconnected, page of it.
\end{beauty}

\section{The \texorpdfstring{$p$}{p}-adic numbers as a locally compact field}\label{sec:padic}

Volume~13 constructs $\Q_p$, the completion of $\Q$ for $|x|_p=p^{-v_p(x)}$, and its ring of integers $\Z_p=\{|x|_p\le1\}$, and proves that $\Z_p$ is compact. The balls $a+p^k\Z_p$ ($a\in\Q_p$, $k\in\Z$) are compact and open, two balls are either disjoint or nested, and
$\Q_p=\bigcup_kp^{-k}\Z_p$ is locally compact and $\sigma$-compact. Those are the facts this section uses. Haar measure is the new ingredient, and the factor by which multiplication scales it is the absolute value. Figure~\ref{fig:zp} shows $\Z_2$ as a tree of balls.

\begin{figure}[ht]
\centering
\fig{ch08-zp}
\caption{The $2$-adic integers as a tree. Each ball $a+2^k\Z_2$ splits into two balls of half the size, according to the next binary digit. A $2$-adic integer
is an infinite path down the tree, and Haar measure gives each ball at depth $k$ the mass $2^{-k}$: it is the law of a sequence of fair coin tosses.}
\label{fig:zp}
\end{figure}

\begin{proposition}\label{prop:haarqp}
There is a unique Haar measure $\mu_p$ on $\Q_p$ with $\mu_p(\Z_p)=1$. It gives each ball $a+p^k\Z_p$ the mass $p^{-k}$, and for $a\in\Q_p^*$ and Borel $E$,
\[\mu_p(aE)=|a|_p\,\mu_p(E).\]
Likewise $\lambda(aE)=|a|\lambda(E)$ for Lebesgue measure on $\R$.
\end{proposition}
\begin{proof}
$\Z_p$ is a compact abelian group, metrizable, so by Corollary~\ref{cor:abelianamenable}(b) it has a unique Haar probability $\mu_0$. The $p^k$ cosets of the subgroup $p^k\Z_p$
in $\Z_p$ are translates of each other and partition $\Z_p$, so each has mass $p^{-k}$. The quotient $\Q_p/\Z_p$ is countable (each coset contains an element of
$\Z[\frac1p]$, the finite sum of the negative powers in the expansion), so we may define $\mu_p(E)=\sum_c\mu_0\bigl((E-c)\cap\Z_p\bigr)$, summing over a set of
representatives $c$ of the cosets. This is a measure, finite on compact sets (which meet finitely many cosets), positive on balls, and invariant, since a
translation permutes the cosets. Uniqueness follows from Theorem~\ref{thm:haaruniq}, and the mass of a ball $a+p^k\Z_p$ with $k<0$ is $p^{-k}$ because it is the union of
$p^{-k}$ translates of $\Z_p$.

For $a\ne0$, $E\mapsto\mu_p(aE)$ is again a Haar measure, since $x\mapsto ax$ is an automorphism of the additive group; so it equals $c(a)\mu_p$, and with $E=\Z_p$,
$c(a)=\mu_p(a\Z_p)=\mu_p(p^{v_p(a)}\Z_p)=p^{-v_p(a)}=|a|_p$. The real case is Volume~7, Theorem~9.7, with $d=1$.
\end{proof}

So the $p$-adic absolute value is not an arbitrary choice: it is the factor by which multiplication by $a$ scales the Haar measure, exactly as $|a|$ does on
$\R$. Write $|\cdot|_\infty=|\cdot|$ for the real absolute value. Then for every $a\in\Q^*$, written $a=\pm\prod_pp^{v_p(a)}$,
\[\prod_{v\in\{\infty,2,3,5,\dots\}}|a|_v=\prod_pp^{v_p(a)}\cdot\prod_pp^{-v_p(a)}=1,\]
the \emph{product formula}: a rational number cannot be large everywhere.

\subsection*{Characters and self-duality}

For $x\in\Q_p$ with expansion $x=\sum_{i\ge-m}a_ip^i$, its \emph{fractional part} is $\{x\}_p=\sum_{i<0}a_ip^i\in\Z[\frac1p]\cap[0,1[$. Then $x-\{x\}_p\in\Z_p$, and
$\{x+y\}_p-\{x\}_p-\{y\}_p\in\Z[\frac1p]\cap\Z_p=\Z$. So
\[\psi_p(x)=e\bigl(\{x\}_p\bigr)\]
is a character of $\Q_p$, continuous because it is constant on the cosets of $\Z_p$. Its kernel is exactly $\Z_p$: $\{x\}_p\in\Z$ only if $\{x\}_p=0$. For $a\in\Q_p$ write
$\psi_{p,a}(x)=\psi_p(ax)$.

\begin{lemma}\label{lem:psiperp}
For $a\in\Q_p$ and $k\in\Z$, $\psi_{p,a}$ is trivial on $p^k\Z_p$ if and only if $a\in p^{-k}\Z_p$.
\end{lemma}
\begin{proof}
$\psi_p(ap^k\Z_p)=1$ means $ap^k\Z_p\subset\Z_p$, that is $|a|_pp^{-k}\le1$.
\end{proof}

\begin{theorem}[$\Q_p$ is self-dual]\label{thm:qpselfdual}
Every character of $\Q_p$ is $\psi_{p,a}$ for a unique $a\in\Q_p$, and $a\mapsto\psi_{p,a}$ is an isomorphism $\Q_p\cong\widehat{\Q_p}$.
\end{theorem}
\begin{proof}
The map is a homomorphism, and it is injective by the lemma: if $\psi_{p,a}=1$, then $a\in p^{-k}\Z_p$ for every $k$, so $a=0$. Let $\chi$ be a character. The
neighbourhood $\{\Re z>0\}$ of $1$ in $\mathbb T$ contains no subgroup except $\{1\}$ (the powers of $z\ne1$ leave it), and by continuity $\chi(p^k\Z_p)$ lies in it for $k$ large; this
image is a subgroup, so $\chi$ is trivial on $p^k\Z_p$. For $n\ge-k$, the characters of the group $p^{-n}\Z_p$ trivial on $p^k\Z_p$ are the characters of the cyclic
group $p^{-n}\Z_p/p^k\Z_p$ of order $p^{n+k}$; there are $p^{n+k}$ of them. By the lemma, the restrictions of the $\psi_{p,a}$ with $a\in p^{-k}\Z_p$ are such characters, and two
of them coincide exactly when $a-a'\in p^n\Z_p$, which leaves $p^{n+k}$ distinct ones. So $\chi=\psi_{p,a_n}$ on $p^{-n}\Z_p$, for some $a_n$ unique modulo $p^n\Z_p$. Then
$a_{n+1}\equiv a_n\pmod{p^n}$, the sequence $(a_n)$ is Cauchy, and its limit $a$ satisfies $a\equiv a_n\pmod{p^n}$, so $\chi=\psi_{p,a}$ on every $p^{-n}\Z_p$, hence on $\Q_p$.
\end{proof}

The Fourier transform on $\Q_p$ is
\[\hat f(y)=\int_{\Q_p}f(x)\overline{\psi_p(xy)}\dd\mu_p(x).\]
The natural test functions are the \emph{Schwartz--Bruhat functions}: locally constant with compact
support, that is, finite linear combinations of indicators of balls; they form a space $\mathcal S(\Q_p)$.

\begin{theorem}[Fourier analysis on $\Q_p$]\label{thm:qpfourier}
\begin{enumerate}
\item[(a)] $\widehat{\mathbf 1_{\Z_p}}=\mathbf 1_{\Z_p}$; more generally $\widehat{\mathbf 1_{a+p^k\Z_p}}(y)=p^{-k}\,\overline{\psi_p(ay)}\,\mathbf 1_{p^{-k}\Z_p}(y)$.
\item[(b)] The transform maps $\mathcal S(\Q_p)$ onto itself, and $\hat{\hat f}(x)=f(-x)$.
\end{enumerate}
\end{theorem}
\begin{proof}
(a) Substituting $x=a+p^kt$ (Proposition~\ref{prop:haarqp}: $\dd x=p^{-k}\dd t$),
\[\widehat{\mathbf 1_{a+p^k\Z_p}}(y)=p^{-k}\,\overline{\psi_p(ay)}\int_{\Z_p}\overline{\psi_p(p^kty)}\dd t.\]
The integrand is a character of the compact group $\Z_p$. If it is trivial, that is $y\in p^{-k}\Z_p$ (Lemma~\ref{lem:psiperp}), the integral is $1$. Otherwise it is $0$: if
$\eta(t_0)\ne1$, invariance gives $\int\eta=\int\eta(t+t_0)\dd t=\eta(t_0)\int\eta$.

(b) By (a) and linearity, $\hat f\in\mathcal S(\Q_p)$ for $f\in\mathcal S(\Q_p)$. Apply (a) twice to $f=\mathbf 1_{a+p^k\Z_p}$: $\hat f=p^{-k}g$ with $g(y)=\overline{\psi_p(ay)}\mathbf 1_{p^{-k}\Z_p}(y)$, and
\[\hat g(x)=\int_{p^{-k}\Z_p}\overline{\psi_p(y(x+a))}\dd y=p^k\,\mathbf 1_{p^k\Z_p}(x+a),\]
by the same computation as in (a) (the integral of a character over the compact group $p^{-k}\Z_p$, of mass $p^k$, trivial exactly when $x+a\in p^k\Z_p$). So
$\hat{\hat f}(x)=\mathbf 1_{p^k\Z_p}(x+a)=\mathbf 1_{-a+p^k\Z_p}(x)=f(-x)$. Since $\hat{\hat f}(-x)=f(x)$, the transform is onto.
\end{proof}

The indicator of $\Z_p$ is to $\Q_p$ what the Gaussian $e^{-\pi x^2}$ is to $\R$ (Chapter~18, equation~\eqref{eq:gaussft}): its own Fourier transform. With $\psi_\infty(x)=e(-x)$ on $\R$, the
real transform $\int f(x)\overline{\psi_\infty(xy)}\dd x$ is the transform of Chapter~18 at $-y$. The minus sign is the one Volume~13 uses, so that a single character on the adèles is trivial on $\Q$.

\begin{beauty}
Where this leads. Theorem~\ref{thm:qpselfdual} is the duality this chapter contributes. Volume~13 constructs $\Q_p$, assembles $\R$ and the fields $\Q_p$ into the adele ring, proves that the diagonal copy of $\Q$ is discrete and that the quotient is compact, and proves the Poisson formula that sums a Schwartz function over $\Q$ on both sides. Applied to a Gaussian at the real place and to the indicator of $\Z_p$ at every prime, that formula is the functional equation of $\zeta(s)$. The local integral that feeds the Euler factor is Exercise~\ref{ex:localzeta}.
\end{beauty}

\section{Exercises}

\begin{exercise}[$\star$]\label{ex:axb}
\emph{The $ax+b$ group.} Let $G=\{(a,b):a>0,b\in\R\}$ with $(a,b)(a',b')=(aa',ab'+b)$, the composition of $x\mapsto a'x+b'$ and then $x\mapsto ax+b$. Show that $\frac{\dd a\dd b}{a^2}$ is a left
Haar measure and $\frac{\dd a\dd b}a$ a right Haar measure, and compute the modular function.
\end{exercise}

\begin{exercise}[$\star\star$]\label{ex:squareszp}
\emph{How many $p$-adic integers are squares?} Let $p$ be odd. Using the criterion for squares in $\Q_p$ from Volume~13, compute $\mu_p(\{x\in\Z_p:x\text{ is a square in }\Q_p\})$. What is the answer for $p=2$?
\end{exercise}

\begin{exercise}[$\star\star$]\label{ex:localzeta}
\emph{The local zeta integral.} Show that for real $s>0$,
\[\int_{\Z_p\setminus\{0\}}|x|_p^{s-1}\dd\mu_p(x)=\frac{1-p^{-1}}{1-p^{-s}},\]
and deduce that $\prod_p\frac1{1-p^{-1}}\int_{\Z_p}|x|_p^{s-1}\dd\mu_p=\zeta(s)$ for $s>1$. Compute also $\int_\R e^{-\pi x^2}|x|^{s-1}\dd x$.
\end{exercise}

\begin{exercise}[$\star\star$]\label{ex:folner}
\emph{F{\o}lner's criterion.} A sequence of finite non-empty sets $F_n\subset G$ is a \emph{F{\o}lner sequence} if $\frac{|gF_n\,\triangle\,F_n|}{|F_n|}\to0$ for every $g\in G$. (a) Show that
$\Z^d$ has one. (b) Show that a group with a F{\o}lner sequence is amenable (take a weak* limit point of the averages over $F_n$). (c) Deduce that $F_2$ has no F{\o}lner
sequence, and check directly that for the ball $B_n$ of radius $n$ in the Cayley graph, $\frac{|aB_n\triangle B_n|}{|B_n|}$ does not tend to $0$.
\end{exercise}

\begin{exercise}[$\star\star$]\label{ex:qzdual}
\emph{A discrete group and its compact dual.} For $a=(a_p)\in\hat{\Z}$ and $x\in\Q$, let $\chi_a(x)=e\bigl(\sum_p\{a_px\}_p\bigr)$. Show that $\chi_a$ is a character of the discrete group $\Q/\Z$, and that
$a\mapsto\chi_a$ is an isomorphism $\hat{\Z}\cong\widehat{\Q/\Z}$.
\end{exercise}

\begin{exercise}[$\star\star$]\label{ex:circlerotation}
\emph{Rotations of the circle.} Let $\alpha$ be irrational and $T$ the rotation $z\mapsto e(\alpha)z$ of $\mathbb T$. (a) Show that the only Borel probability on $\mathbb T$ invariant under $T$ is
the Haar measure (compare Fourier coefficients). (b) Show, using Theorem~\ref{thm:banachmeasure}'s method with the amenable group $\Z$, that there is a finitely additive
probability on \emph{all} subsets of $\mathbb T$ invariant under $T$. (c) Show that such a probability cannot be countably additive (use a set of representatives of
$\mathbb T/\langle e(\alpha)\rangle$, as in Vitali's construction).
\end{exercise}

\section{Hints}

\begin{hint}{ex:axb}
Left multiplication by $(a_0,b_0)$ is $(a,b)\mapsto(a_0a,a_0b+b_0)$, with Jacobian $a_0^2$; right multiplication is $(a,b)\mapsto(aa_0,ab_0+b)$, with Jacobian $a_0$.
\end{hint}
\begin{hint}{ex:squareszp}
Split $\Z_p\setminus\{0\}$ into the shells $p^v\Z_p^*$, of mass $p^{-v}(1-\frac1p)$.
\end{hint}
\begin{hint}{ex:localzeta}
On the shell $p^v\Z_p^*$, $|x|_p=p^{-v}$. For the real integral, substitute $u=\pi x^2$.
\end{hint}
\begin{hint}{ex:folner}
(b) Let $m_n(f)=\frac1{|F_n|}\sum_{x\in F_n}f(x)$ and estimate $|m_n(\lambda_gf)-m_n(f)|$. (c) $aB_n\setminus B_n$ contains all words of length $n+1$ beginning with $a$.
\end{hint}
\begin{hint}{ex:qzdual}
A character of $\Q/\Z$ is determined by its values on $\frac1{n}$, $n\ge1$; restrict to the $p$-primary part $\Z[\frac1p]/\Z$ and use Theorem~\ref{thm:qpselfdual} for the characters of
$\Q_p$ trivial on $\Z_p$.
\end{hint}
\begin{hint}{ex:circlerotation}
(a) The $k$-th Fourier coefficient of an invariant measure is multiplied by $e(-k\alpha)$. (b) Average a non-invariant extension over $n\in\Z$ with an invariant mean.
\end{hint}

\section{Solutions}

\begin{solution}{ex:axb}
\idea Haar measures are found by asking which density is preserved by the Jacobian of the translations; left and right translations have different
Jacobians in a non-commutative group.
\writeup Left multiplication $L(a,b)=(a_0a,a_0b+b_0)$ is affine with Jacobian determinant $a_0^2$. For a density $h$, invariance means $h(L(a,b))a_0^2=h(a,b)$; $h=a^{-2}$
works: $(a_0a)^{-2}a_0^2=a^{-2}$. By the change of variables formula, $\int_{L(E)}a^{-2}\dd a\dd b=\int_Ea^{-2}\dd a\dd b$. Right multiplication $R(a,b)=(aa_0,ab_0+b)$ has
Jacobian matrix $\begin{pmatrix}a_0&0\\b_0&1\end{pmatrix}$, determinant $a_0$, and $h=a^{-1}$ works. With $\mu=a^{-2}\dd a\dd b$, $\mu(Eg_0)=\int_Ea^{-2}a_0^{-2}\cdot a_0\dd a\dd b$ for
$g_0=(a_0,b_0)$, since $R$ multiplies $a$ by $a_0$ and has Jacobian $a_0$. So $\Delta(a_0,b_0)=a_0^{-1}$, and the group is not unimodular. (Conventions differ: some books
define $\Delta$ with $g^{-1}$, which gives $a_0$.)
\end{solution}

\begin{solution}{ex:squareszp}
\idea A $p$-adic number is a square according to the parity of its valuation and the residue of its unit part. Haar measure makes valuation and unit part
independent, and the unit part uniform modulo $p$.
\writeup The shell $p^v\Z_p^*=p^v\Z_p\setminus p^{v+1}\Z_p$ has mass $p^{-v}-p^{-v-1}$. On it, by the criterion for squares in $\Q_p$ (Volume~13), $x=p^vu$ is a square if and only if $v$ is even and
$u\bmod p$ is a non-zero square; the $p-1$ unit residues are equally likely ($\Z_p^*$ is the union of $p-1$ cosets of $p\Z_p$, all of mass $\frac1p$) and half of them
are squares. So the measure is
\[\sum_{v\text{ even}}\frac12\bigl(p^{-v}-p^{-v-1}\bigr)=\frac12\Bigl(1-\frac1p\Bigr)\frac1{1-p^{-2}}=\frac p{2(p+1)}.\]
For $p=3$ this is $\frac38$. For $p=2$, a unit $u$ is a square if and only if $u\equiv1\pmod8$, which is one of the four unit classes modulo $8$, so the shell $2^v\Z_2^*$ contributes
$\frac14\cdot2^{-v-1}$ for even $v$, and the total is $\frac18\cdot\frac1{1-\frac14}=\frac16$.
\end{solution}

\begin{solution}{ex:localzeta}
\idea Both integrals are sums over the natural ``shells'' of the absolute value: geometric series at $p$, and a Gamma integral at infinity. Their
product over all places is the completed zeta function.
\writeup The point $0$ has measure $0$. On the shell $p^v\Z_p^*$, of mass $p^{-v}(1-p^{-1})$, the integrand is $p^{-v(s-1)}$, so the integral is
\[\sum_{v\ge0}p^{-v(s-1)}p^{-v}(1-p^{-1})=(1-p^{-1})\sum_{v\ge0}p^{-vs}=\frac{1-p^{-1}}{1-p^{-s}}.\]
Dividing by $1-p^{-1}$ leaves the Euler factor $(1-p^{-s})^{-1}$, and the Euler product (Volume~6) gives $\zeta(s)$ for $s>1$. At infinity, by parity and $u=\pi x^2$
($\dd x=\frac12\pi^{-1/2}u^{-1/2}\dd u$),
\[\int_\R e^{-\pi x^2}|x|^{s-1}\dd x=2\int_0^\infty e^{-u}\Bigl(\frac u\pi\Bigr)^{\frac{s-1}2}\frac{\dd u}{2\sqrt{\pi u}}=\pi^{-s/2}\Gamma\Bigl(\frac s2\Bigr).\]
The product of all the local factors is Riemann's completed zeta function $\pi^{-s/2}\Gamma(\frac s2)\zeta(s)$, the function that is symmetric under $s\mapsto1-s$. Volume~13 explains the symmetry by the Poisson formula on the adèles.
\end{solution}

\begin{solution}{ex:folner}
\idea Averages over sets that barely move under translation are almost invariant; a limit of almost invariant means is invariant. Free groups are
too large: a ball is dominated by its boundary.
\writeup (a) $F_n=\{0,\dots,n-1\}^d$: $|(g+F_n)\triangle F_n|\le2d|g|_\infty n^{d-1}$, so the ratio is $O(\frac1n)$. (b) Since $m_n(\lambda_gf)=\frac1{|F_n|}\sum_{x\in F_n}f(g^{-1}x)=\frac1{|F_n|}\sum_{y\in g^{-1}F_n}f(y)$, we get
$|m_n(\lambda_gf)-m_n(f)|\le\frac{|g^{-1}F_n\triangle F_n|}{|F_n|}\|f\|_\infty\to0$. By Banach--Alaoglu the $m_n$ have a weak* limit point $m$ (a limit of a subnet), which is a mean, and for each $f$ and $g$,
$m(\lambda_gf)-m(f)$ is a limit point of $m_n(\lambda_gf)-m_n(f)\to0$, so it is $0$. (c) By (b) and Proposition~\ref{prop:f2paradox}, $F_2$ has no F{\o}lner sequence. Directly: the
sphere of radius $k\ge1$ has $4\cdot3^{k-1}$ words, so $|B_n|=2\cdot3^n-1$, and the $3^n$ words of length $n+1$ beginning with $a$ are $a$ times words of length $n$ not beginning
with $a^{-1}$, so they lie in $aB_n\setminus B_n$. Hence $\frac{|aB_n\triangle B_n|}{|B_n|}\ge\frac{3^n}{2\cdot3^n}=\frac12$.
\end{solution}

\begin{solution}{ex:qzdual}
\idea $\Q/\Z$ is the direct sum of its $p$-parts $\Z[\frac1p]/\Z\cong\Q_p/\Z_p$, and a character of $\Q_p$ trivial on $\Z_p$ comes from a $p$-adic integer by
Theorem~\ref{thm:qpselfdual}. Putting the primes together gives $\prod_p\Z_p=\hat{\Z}$.
\writeup The sum is finite, since $a_px\in\Z_p$ for $p$ not dividing the denominator of $x$. As in Section~\ref{sec:padic}, $\{u+v\}_p-\{u\}_p-\{v\}_p\in\Z$, so $\chi_a$ is a homomorphism
$\Q\to\mathbb T$, and $\chi_a(n)=1$ for $n\in\Z$ ($a_pn\in\Z_p$); so it is a character of $\Q/\Z$ (every homomorphism is continuous on a discrete group). Clearly $\chi_{a+b}=\chi_a\chi_b$.
For $x=p^{-k}$, the terms with $q\ne p$ vanish ($a_qp^{-k}\in\Z_q$), so $\chi_a(p^{-k})=\psi_p(a_pp^{-k})$.

Injectivity: if $\chi_a=1$, then $\psi_{p,a_p}$ is trivial on $p^{-k}\Z$ and on $\Z_p$, hence on $p^{-k}\Z+\Z_p=p^{-k}\Z_p$, for every $k$; by Lemma~\ref{lem:psiperp}, $a_p\in p^k\Z_p$ for all $k$, so
$a_p=0$. Surjectivity: let $\chi$ be a character of $\Q/\Z$. Since $\Z[\frac1p]+\Z_p=\Q_p$ and $\Z[\frac1p]\cap\Z_p=\Z$, the group $\Z[\frac1p]/\Z$ is isomorphic to $\Q_p/\Z_p$, and the restriction of $\chi$
becomes a character of $\Q_p$ trivial on the open subgroup $\Z_p$ (hence continuous). By Theorem~\ref{thm:qpselfdual} and Lemma~\ref{lem:psiperp} with $k=0$, it is $\psi_{p,a_p}$ with
$a_p\in\Z_p$. Then $\chi$ and $\chi_a$, $a=(a_p)$, agree on every $\Z[\frac1p]/\Z$, and these subgroups generate $\Q/\Z$ (partial fractions: every rational is an integer plus a sum of
fractions with prime power denominators), so $\chi=\chi_a$.
\end{solution}

\begin{solution}{ex:circlerotation}
\idea Unique ergodicity is a statement about Fourier coefficients, while invariant finitely additive measures on all sets come from amenability. The two
together show how far finitely additive measures are from countably additive ones.
\writeup (a) If $\nu$ is $T$-invariant, its coefficients $c_k=\int z^{-k}\dd\nu$ satisfy $c_k=\int(e(\alpha)z)^{-k}\dd\nu=e(-k\alpha)c_k$, so $c_k=0$ for $k\ne0$ ($k\alpha\notin\Z$). Trigonometric
polynomials are dense in $C(\mathbb T)$ (Fej\'er, or Stone--Weierstrass), so $\int f\dd\nu=\int f\dd\mu$ for all continuous $f$, and $\nu=\mu$ by the end of the proof of
Theorem~\ref{thm:haaruniq}. (b) As in Theorem~\ref{thm:banachmeasure}, Hahn--Banach extends $f\mapsto\int f\dd\mu$ to a positive $L$ on all bounded functions on $\mathbb T$ with $L(\mathbf 1)=1$
(take $p$ the upper integral). With an invariant mean $m$ on $\Z$ (Corollary~\ref{cor:abelianamenable}), $\nu(A)=m\bigl(n\mapsto L(\mathbf 1_{T^{-n}A})\bigr)$ is finitely additive, $T$-invariant,
and $\nu(\mathbb T)=1$. (c) Let $V$ contain one point of each orbit of the group $\langle e(\alpha)\rangle$ acting by rotation. The sets $T^nV$, $n\in\Z$, are disjoint and cover
$\mathbb T$, and they all have the same $\nu$-measure. If $\nu$ were countably additive, $1=\sum_n\nu(T^nV)$ would be $0$ or $\infty$.
\end{solution}

\section*{Sources and further reading}
\addcontentsline{toc}{section}{Sources and further reading}

A.~Haar, ``Der Massbegriff in der Theorie der kontinuierlichen Gruppen'', \emph{Ann.\ of Math.} 34 (1933), 147--169; A.~Weil, \emph{L'int\'egration dans les groupes topologiques
et ses applications} (Hermann, 1940). The fixed point theorem is in A.~Markov, \emph{Dokl.\ Akad.\ Nauk SSSR} 10 (1936), 311--314, and S.~Kakutani, \emph{Proc.\ Imp.\ Acad.\ Tokyo} 14
(1938), 242--245. Amenability begins with J.~von Neumann, ``Zur allgemeinen Theorie des Masses'', \emph{Fund.\ Math.} 13 (1929), 73--116; the name is M.~M.~Day's, \emph{Illinois J.\
Math.} 1 (1957), 509--544. Banach's measure is in S.~Banach, ``Sur le probl\`eme de la mesure'', \emph{Fund.\ Math.} 4 (1923), 7--33, the paradox in S.~Banach and A.~Tarski, \emph{Fund.\
Math.} 6 (1924), 244--277, and Hausdorff's in his \emph{Grundz\"uge der Mengenlehre} (1914). The best book on all of this is G.~Tomkowicz and S.~Wagon, \emph{The Banach--Tarski
Paradox}, 2nd ed. (Cambridge, 2016).

L.~Pontryagin, ``The theory of topological commutative groups'', \emph{Ann.\ of Math.} 35 (1934), 361--388; E.~R.~van Kampen, ``Locally bicompact abelian groups and their character
groups'', \emph{Ann.\ of Math.} 36 (1935), 448--463. For the general theory: G.~B.~Folland, \emph{A Course in Abstract Harmonic Analysis}, 2nd ed. (CRC, 2016). J.~T.~Tate's thesis,
``Fourier analysis in number fields and Hecke's zeta-functions'' (Princeton, 1950), is printed in J.~W.~S.~Cassels and A.~Fr\"ohlich (eds.), \emph{Algebraic Number Theory}
(Academic Press, 1967), 305--347, and is beautifully readable; D.~Ramakrishnan and R.~J.~Valenza, \emph{Fourier Analysis on Number Fields} (Springer, 1999), develops everything
needed for it.
